\documentclass[11pt]{article}
\usepackage[T1]{fontenc}
\usepackage{amsmath,amssymb,amsthm}
\usepackage[margin=1in]{geometry}
\usepackage{booktabs}
\usepackage[hidelinks]{hyperref}

\newtheorem{theorem}{Theorem}
\newtheorem{lemma}[theorem]{Lemma}
\newcommand{\TR}{\mathcal T_{\mathbb R}}
\newcommand{\TZ}{\mathcal T_{\mathbb Z}}

\title{The Exact Tur\'an Constant for Three Equal Intervals at
  $\lambda=5/2$}
\author{Xiaojian Zeng\\
PBC School of Finance, Tsinghua University, 100190\\
\texttt{xiaojianzeng@foxmail.com}}
\date{}

\begin{document}
\maketitle

\begin{abstract}
Let
\[
 \Omega_{5/2}=(-1,1)\cup(3/2,7/2)\cup(-7/2,-3/2).
\]
We determine its Tur\'an constant exactly.  If $u$ is the unique root in
$[533/1000,267/500]$ of
\[
256u^7+256u^6-96u^5-80u^4+40u^3+30u^2-15u-5=0,
\]
then
\[
 \mathcal T_{\mathbb R}(\Omega_{5/2})
 =\frac{(2u+1)^2(4u^2-6u+1)^2}{128u^6-20u^2+5}
 =2.1353071788763974322\ldots .
\]
The lower certificate is the normalized square of a cubic in the Chebyshev
basis.  A matching upper certificate is a positive three-atomic moment
functional supported at the cubic's roots.  The result determines the value
but does not classify all extremizers.
\end{abstract}

\section{Reduction and statement}

For an open symmetric set $\Omega\subset\mathbb R$, its Tur\'an constant is
the supremum of $\int_{\mathbb R}f$ over continuous positive-definite
functions $f$ satisfying $f(0)=1$ and $\operatorname{supp}f\subset\Omega$.
Kolountzakis and R\'ev\'esz studied this question on arbitrary locally compact
abelian groups and for more general domains
\cite{KolountzakisRevesz2006}.  Fu, Li, Wang, and Wang recently treated the
three-interval family
\[
 \Omega_\lambda=(-1,1)\cup(\lambda-1,\lambda+1)
 \cup(-\lambda-1,-\lambda+1).
\]
Their rational-parameter reduction \cite[Theorem~1.2]{FuLiWangWang2026}
gives, at $\lambda=5/2$,
\begin{equation}\label{eq:reduction}
 \TR(\Omega_{5/2})=\frac12\TZ(A),
 \qquad A=\{0,\pm1,\pm4,\pm5,\pm6\}.
\end{equation}
Equivalently, writing $T_k(\cos t)=\cos(kt)$, $\TZ(A)$ is the maximum
$M$ of
\begin{equation}\label{eq:finite}
 p_c(1)=1+2(c_1+c_4+c_5+c_6)
\end{equation}
over all
\[
 p_c(x)=1+2c_1T_1(x)+2c_4T_4(x)+2c_5T_5(x)+2c_6T_6(x)
\]
that are nonnegative on $[-1,1]$.  The same source proved
\[
 \frac{559+80\sqrt{42}}{506}
 \leq \TR(\Omega_{5/2})\leq\frac94
\]
and explicitly did not claim optimality for its finite feasible polynomial
\cite[Proposition~4.2 and Lemma~5.1]{FuLiWangWang2026}.

Put
\begin{equation}\label{eq:P}
 P(U)=256U^7+256U^6-96U^5-80U^4+40U^3+30U^2-15U-5.
\end{equation}

\begin{theorem}\label{thm:main}
Let $u$ be the unique zero of $P$ in
$I=[533/1000,267/500]$.  Then
\begin{equation}\label{eq:value}
 \TR(\Omega_{5/2})
 =\frac{(2u+1)^2(4u^2-6u+1)^2}
 {128u^6-20u^2+5}.
\end{equation}
\end{theorem}

We prove the theorem by displaying both sides of an exact finite
primal--dual certificate.  Zeng, Liu, and Ratnavelu give exact
sum-of-squares and positive-semidefinite Gram certificates for another sharp
finite-dimensional inequality \cite{ZengLiuRatnavelu2026}.  The present dual
certificate instead takes the form of an exact atomic moment functional.

\section{The exact certificate}

First, the algebraic parameter is unambiguous.  Direct evaluation gives
\[
 P(533/1000)=-\frac{14387903930821523}{3906250000000000000}<0,
 \qquad
 P(267/500)=\frac{2066305610212673}{30517578125000000}>0.
\]
Moreover, throughout $I$, termwise rational interval bounds in $P'$ give
\[
 P'(U)\geq
 \frac{275284068057136783}{3906250000000000}>0.
\]
Thus $P$ has exactly one zero $u$ in $I$.

\subsection{Primal certificate}

Set
\begin{equation}\label{eq:qA}
 q(X)=X^3+uX^2+(2u^2-5/4)X-2u^3,
 \qquad A_0=\frac{128u^6-20u^2+5}{32}.
\end{equation}
An exact change from the monomial basis to the Chebyshev basis gives
\begin{align}\label{eq:cheb-square}
q(X)^2={}&A_0T_0(X)
-\frac{u(64u^4-40u^2+5)}8T_1(X)
+\frac{5u^2-1}{8}T_4(X)\notag\\
&+\frac u8T_5(X)+\frac1{32}T_6(X).
\end{align}
In particular, the $T_2$ and $T_3$ coefficients vanish identically.
Also $A_0>0$, since it is the constant Chebyshev coefficient, equivalently
the arcsine mean, of the nonzero square $q^2$.  Hence
\begin{equation}\label{eq:primal}
 p(X)=\frac{q(X)^2}{A_0}
\end{equation}
is feasible for \eqref{eq:finite}.  Since
\[
 q(1)=-\frac{(2u+1)(4u^2-6u+1)}4,
\]
this gives
\begin{equation}\label{eq:primal-value}
 \frac{p(1)}2=
 \frac{(2u+1)^2(4u^2-6u+1)^2}{128u^6-20u^2+5}.
\end{equation}

\subsection{Positive atomic dual certificate}

Define
\[
 D=5(2u-1)(32u^4+8u^2+1)
\]
and moments
\begin{align*}
 m_0&=-\frac{512u^4-256u^3-32u^2+6u-3}{D},
 &m_1&=-1,\\
 m_2&=-\frac{256u^6-128u^5+48u^4-44u^3+42u^2-23u-1}{D}.
\end{align*}
For $0\leq n\leq3$, recursively put
\begin{equation}\label{eq:recurrence}
 m_{n+3}=-um_{n+2}-(2u^2-5/4)m_{n+1}+2u^3m_n,
\end{equation}
and define $L(X^j)=m_j$ for $0\leq j\leq6$.
Direct reduction modulo $P(u)=0$ gives
\begin{equation}\label{eq:dual-moments}
 L(T_1)=L(T_4)=L(T_5)=L(T_6)=-1.
\end{equation}
For transparency, the $T_1,T_4,T_6$ identities simplify directly, while
the residual in the $T_5$ identity is
\[
 -\frac{4(4u^2-6u+1)P(u)}
 {5(2u-1)(32u^4+8u^2+1)}=0.
\]

It remains to show that $L$ is positive on polynomials nonnegative on
$[-1,1]$.  The following rational interval enclosures hold uniformly for
$U\in I$; each is obtained by substituting the interval $I$ into
$q_U(x)=x^3+Ux^2+(2U^2-5/4)x-2U^3$ with outward rational endpoints.
\begin{center}
\small
\begin{tabular}{c@{\qquad}c}
\toprule
$x$ & enclosure for $q_U(x)$ \\
\midrule
$-23/25$ & $[-424403/62500000,-1136517/500000000]$\\
$-91/100$ & $[443943/250000000,3126773/500000000]$\\
$-21/50$ & $[1671/1953125,1817983/500000000]$\\
$-41/100$ & $[-1299557/250000000,-1223727/500000000]$\\
$79/100$ & $[-2187711/125000000,-6741997/500000000]$\\
$4/5$ & $[194737/62500000,3585363/500000000]$\\
\bottomrule
\end{tabular}
\end{center}
The alternating signs show that $q$ has one root in each of
\begin{equation}\label{eq:root-intervals}
 (-23/25,-91/100),\qquad(-21/50,-41/100),\qquad(79/100,4/5).
\end{equation}
These three roots are distinct and lie in $(-1,1)$.

Let $H=(m_{i+j})_{0\leq i,j\leq2}$.  Its leading principal minors are
\begin{align}
\Delta_1={}&\frac{-512u^4+256u^3+32u^2-6u+3}
 {5(2u-1)(32u^4+8u^2+1)},\label{eq:d1}\\
\Delta_2={}&\frac{(4u^2-6u+1)
 (64u^4-32u^3-24u^2+32u-11)}
 {25(2u-1)^2(32u^4+8u^2+1)^2}\notag\\
&\quad\cdot(112u^4+112u^3+46u^2+3u+2),\label{eq:d2}\\
\Delta_3={}&\frac{(2u+1)(4u^2-6u+1)^2
 (64u^4-32u^3-24u^2+32u-11)^2}
 {2000(2u-1)^3(32u^4+8u^2+1)^3}\notag\\
&\quad\cdot(512u^5+1024u^4+864u^3+328u^2+32u-11).
\label{eq:d3}
\end{align}
Here all denominators are positive.  For completeness, outward Horner
evaluation on $I$ gives the following exact enclosures for the remaining
sign factors, in their order of appearance:
\newpage
\begin{center}
\footnotesize
\begin{tabular}{c@{\enspace}c}
\toprule
factor & enclosure on $I$\\
\midrule
$-512U^4+256U^3+32U^2-6U+3$
 & $[6105461933/976562500,1551025452/244140625]$\\
$4U^2-6U+1$
 & $[-133189/125000,-132439/125000]$\\
$64U^4-32U^3-24U^2+32U-11$
 & $[-3551016821/7812500000,-805081621/1953125000]$\\
$112U^4+112U^3+46U^2+3U+2$
 & $[2666575353447/62500000000,167503702397/3906250000]$\\
$512U^5+1024U^4+864U^3+328U^2+32U-11$
 & $[653772656842393/1953125000000,20549356967107/61035156250]$\\
\bottomrule
\end{tabular}
\end{center}
Consequently $\Delta_1,\Delta_2,\Delta_3>0$, so $H$ is positive definite.

The recurrence \eqref{eq:recurrence} says that
$L(qs)=0$ for every polynomial $s$ of degree at most three.  Let
$r_1,r_2,r_3$ be the roots in \eqref{eq:root-intervals}, and let $h_i$ be
the degree-two Lagrange polynomial with $h_i(r_j)=\delta_{ij}$.  Reduction
modulo $q$ and interpolation give, for every polynomial $f$ of degree at
most six,
\begin{equation}\label{eq:atomic}
 L(f)=\sum_{i=1}^3w_i f(r_i),\qquad w_i=L(h_i).
\end{equation}
Since $h_i^2-h_i$ is divisible by $q$, and its quotient has degree one,
\[
 w_i=L(h_i)=L(h_i^2)>0;
\]
the last inequality follows from positive definiteness of $H$.  Thus
\eqref{eq:atomic} is a positive three-atomic measure supported in $(-1,1)$.

For any feasible $p_c$ in \eqref{eq:finite}, equations
\eqref{eq:dual-moments} and \eqref{eq:atomic} yield
\[
 0\leq L(p_c)=m_0-2(c_1+c_4+c_5+c_6),
 \qquad p_c(1)\leq1+m_0.
\]
For the primal polynomial \eqref{eq:primal}, the recurrence gives
$L(q^2)=0$, so equality holds: $p(1)=1+m_0$.  Hence $M=p(1)$.
Combining this equality with \eqref{eq:reduction} and
\eqref{eq:primal-value} proves Theorem~\ref{thm:main}.

\section{Scope and verification}

The argument determines the exact constant in the full continuous problem
through the cited finite reduction.  It does not classify all extremizing
positive-definite functions.  No assertion of priority or of formal
open-problem status is made.

Floating-point linear programming was used only to locate a candidate
contact pattern.  The displayed proof uses exact symbolic identities and
rational interval enclosures.  A deterministic SymPy verifier accompanying
the source recomputes the Chebyshev expansion, moment identities, principal
minors, rational sign enclosures, complementarity, and the reported decimal.
OpenAI Codex assisted with exploratory symbolic computation, verification
code, and language editing.  The author reviewed the mathematical argument
and is responsible for its claims.

\small
\bibliographystyle{alphaurl}
\bibliography{references}

\end{document}
